%%==========================================================================%%
%% datarecords.tex -- "Data Records".                                        %%
%% DICOM release (AutoPET-shaped): DICOM hierarchy + a per-examination   %%
%% metadata CSV + label records. American spelling. Cohort numbers are       %%
%% [TBD] pending reconciliation.                                             %%
%%==========================================================================%%

\section*{Data Records}\label{sec:datarecords}

The dataset is deposited in [TBD: repository] under accession [TBD] and Digital Object Identifier [TBD], under [TBD: license, pending the open-versus-controlled-access decision].
The imaging, segmentation labels, and accompanying metadata are organized as described below.

\subsection*{File organization}\label{subsec:layout}
The imaging follows the standard DICOM hierarchy.
Each subject contains one or more imaging studies (examinations), and each study contains the PET and MRI series acquired in that session.
Studies and series are identified by their DICOM unique identifiers, which are remapped during de-identification, and a per-subject date offset is applied so that the relative timing between a subject's examinations is preserved.

\subsection*{Imaging data}\label{subsec:imaging}
The PET and MRI image data are provided in DICOM, the native clinical format, with the acquisition parameters retained in the DICOM headers.
Segmentation labels are provided as [TBD: NIfTI or DICOM-SEG] volumes spatially aligned to the corresponding imaging series.

\subsection*{Metadata}\label{subsec:metadata}
A per-examination metadata table is provided as an open plain-text file (CSV) with an accompanying data dictionary, so the dataset is usable with standard tooling.
The table captures the indexing variables useful for cohort selection and modeling, grouped as in Table~\ref{tab:vars}.

\begin{table}[h]
\caption{Indexing variables provided with the dataset [TBD: finalize].}\label{tab:vars}
\begin{tabular}{@{}ll@{}}
\toprule
Group & Example fields \\
\midrule
Subject / demographics & subject\_id, age\_at\_scan, sex \\
Clinical / presentation & indication, normal\_vs\_abnormal, pathology\_category, diagnosis \\
Acquisition & scanner, field\_strength, sequence/contrast, voxel\_spacing, TR/TE \\
Labels / annotation & label\_id, structure, label\_type \\
File / QC & dicom\_reference, checksum, qc\_status, license \\
\botrule
\end{tabular}
\end{table}

\subsection*{Label records}\label{subsec:labels}
The released segmentation labels fall into three categories.
Whole-body anatomical segmentation labels are provided for [TBD] examinations.
Tumor-lesion segmentation labels are provided for [TBD] examinations.
A further [TBD] examinations are disease-free follow-up scans demonstrating complete metabolic response, recorded as such so they can be selected as pathology-negative references.
Each label records its target structure and label type, and is stored [TBD: state file formats and folder location, e.g. under a derivatives folder mirroring the imaging layout].
